\scanpage{38}
affine de $X$, $\underline{x}=(x_1,-,x_{\unsure{n|m}})$ est un syst. de générateurs de son anneau affine sur $K$, $c$ un $M_K$-diviseur, et $E(\text{---})$ est l'ens des $(P,v)\in\deleted{X}U(\overline K)\times M_{\overline K}$ tels que
\[
v(\underline{x}(P))=\mathop{\mathrm{Inf}}\limits_{1\leq i\leq n} v(x_i(P))\geq c(v)
\]
$\bigl\{$ i.e. \inserted{essentiellement} il existe une partie finie $S$ de $M_K$ telle que $E$ est formé des $(P,v)$ \inserted{$\in U(\overline K)\times M_{\overline K}$} tels que
\[
\begin{aligned}
&\text{a) }P\text{ soit entier pour }\underline{x},v\quad\text{si}\quad v|K\in M_K-S\\
&\text{b) }v(\unsure{\underline{x}(P)|x(P)})\geq-M\quad\text{si}\quad v|K\in S\quad\bigl\}.
\end{aligned}
\]

% slanted
\margin{\uline{N.B.} La notion de partie bornée est inv. par ext. finie du corps de base. La fonctorialité pour $X\to Y$ n'est pas claire, aussi il y aurait peut-être lieu de remplacer $U\unsure{\hookrightarrow|\subset}X$ par $U\to X$ ----}
% slanted
\margin{\uline{Question.} Dans le cas ``global'', une partie ``$M_K$-bornée'' de $\mathbb{P}^{r}_K$ est-elle \uline{de} \uline{hauteur} \uline{bornée} ?}

\uline{Fonctions admissibles}

Une fonction $\varphi:X(\overline K)\times M_{\overline K}\longrightarrow\mathbb R$ est dite \uline{admissible} si elle est \deleted{\unsure{bor}} $M_K$-bornée sur toute partie $M_K$-bornée. Comme les ens $\{P\}\times M_{\overline K}$, où $P\in X(\overline K)$, sont $M_K$-bornés, on voit que \deleted{si} pour tt $P\in\deleted{X}U(\overline K)_{\unsure{\theta|{}}}$ fixé, la fonction
\[
\widetilde{\varphi}(P)=(v\mapsto\varphi(P,v)):M_{\overline K}\longrightarrow\mathbb R
\]
a la vertu d'être \unsure{majorée|majoré} \inserted{en module} par une fonction de la forme $c(\unsure{w|v})$, avec
% archive sheet number: 37
